\chapter{Peano Axioms and Arithmetic from Successor}

\section{Purpose of the Note}

This independent note explains how the natural numbers can be described without
assuming ordinary arithmetic in advance. The starting point is not a list of
calculation rules, but a small structural idea: there is a first number, called
\(0\), and there is a successor operation which produces the next number.

The main goal is to show how addition and multiplication can be introduced from
this successor structure, and then to verify concrete equalities such as
\(2+3=5\) and \(2\cdot 2=4\) by unfolding the definitions.

\section{The Peano Structure}

\begin{definitionbox}
A \emph{Peano system} consists of a set \(\mathbb{N}\), a distinguished element
\(0\in \mathbb{N}\), and a function
\[
S:\mathbb{N}\to\mathbb{N},
\]
called the \emph{successor function}, satisfying the following axioms.

\begin{enumerate}
    \item \(0\) is a natural number.
    \item If \(n\in\mathbb{N}\), then \(S(n)\in\mathbb{N}\).
    \item \(0\) is not the successor of any natural number: there is no
    \(n\in\mathbb{N}\) such that \(S(n)=0\).
    \item The successor function is injective: if \(S(m)=S(n)\), then
    \(m=n\).
    \item If a subset \(A\subseteq \mathbb{N}\) contains \(0\), and whenever it
    contains \(n\) it also contains \(S(n)\), then \(A=\mathbb{N}\).
\end{enumerate}
\end{definitionbox}

The fifth axiom is the \emph{principle of induction}. It says that a property
which holds for \(0\) and is preserved by the successor operation must hold for
every natural number.

The usual numerals are then abbreviations:
\[
1=S(0),\qquad 2=S(1),\qquad 3=S(2),\qquad 4=S(3),\qquad 5=S(4).
\]
Thus the statement \(2+3=5\) should not be read as something already known. In
this note it will be derived from the recursive definition of addition.

\section{The Recursion Principle}

The Peano axioms support a basic construction principle: to define a function
on \(\mathbb{N}\), it is enough to specify its value at \(0\) and explain how to
pass from the value at \(n\) to the value at \(S(n)\).

\begin{theorembox}
Let \(X\) be a set, let \(x_0\in X\), and let \(F:X\to X\) be a function. There
is a unique function \(f:\mathbb{N}\to X\) such that
\[
f(0)=x_0,
\qquad
f(S(n))=F(f(n))
\quad\text{for every } n\in\mathbb{N}.
\]
\end{theorembox}

In this note, the recursion principle is used to define addition first, and
multiplication second. The definitions are recursive in the second argument.

\section{Addition}

\begin{definitionbox}
For natural numbers \(a,b\in\mathbb{N}\), addition is defined recursively by
\[
a+0=a,
\]
and
\[
a+S(b)=S(a+b).
\]
\end{definitionbox}

The first rule says that adding zero on the right changes nothing. The second
rule says that adding the successor of \(b\) is the same as taking one successor
after adding \(b\).

\begin{proposition}
For every \(b\in\mathbb{N}\), one has
\[
0+b=b.
\]
\end{proposition}

\begin{proof}
We prove the statement by induction on \(b\).

For \(b=0\), the definition of addition gives
\[
0+0=0.
\]

Assume now that \(0+b=b\). Then
\[
0+S(b)=S(0+b)=S(b),
\]
where the first equality is the recursive definition of addition and the second
uses the induction hypothesis. Hence the property is preserved by successor.
By induction, \(0+b=b\) for every \(b\in\mathbb{N}\).
\end{proof}

\subsection*{A Formal Computation of \(2+3\)}

We first compute small additions directly from the recursive rules.

Since \(1=S(0)\),
\[
2+1=2+S(0)=S(2+0)=S(2)=3.
\]
Since \(2=S(1)\),
\[
2+2=2+S(1)=S(2+1)=S(3)=4.
\]
Since \(3=S(2)\),
\[
2+3=2+S(2)=S(2+2)=S(4)=5.
\]
Therefore
\[
\boxed{2+3=5}.
\]

Notice that the computation did not assume a prior addition table. It used only
the definitions of the numerals and the two recursive clauses defining addition.

\section{Multiplication}

\begin{definitionbox}
For natural numbers \(a,b\in\mathbb{N}\), multiplication is defined recursively
by
\[
a\cdot 0=0,
\]
and
\[
a\cdot S(b)=a\cdot b+a.
\]
\end{definitionbox}

The definition says that multiplying by zero gives zero, and multiplying by the
successor of \(b\) means multiplying by \(b\) and then adding one more copy of
\(a\).

\subsection*{A Formal Computation of \(2\cdot 2\)}

Since \(1=S(0)\),
\[
2\cdot 1=2\cdot S(0)=2\cdot 0+2=0+2=2,
\]
where the last equality follows from the proposition \(0+b=b\).

Since \(2=S(1)\),
\[
2\cdot 2=2\cdot S(1)=2\cdot 1+2=2+2=4.
\]
Therefore
\[
\boxed{2\cdot 2=4}.
\]

Again, the equality is not taken from a multiplication table. It follows from
the recursive definition of multiplication together with the previously defined
addition.

\section{What Has Been Constructed}

The order of construction matters.

\begin{enumerate}
    \item The Peano axioms give the natural numbers as a successor structure.
    \item The numerals \(1,2,3,4,5,\ldots\) are names for repeated successors of
    \(0\).
    \item Addition is defined by recursion from the successor operation.
    \item Multiplication is defined by recursion from addition.
    \item Concrete arithmetic facts are then consequences of these definitions.
\end{enumerate}

\begin{summarybox}
The Peano axioms do not start with arithmetic tables. They start with \(0\), the
successor function \(S\), injectivity of successor, the fact that \(0\) has no
predecessor, and induction. From this structure one defines addition, then
multiplication, and only afterward proves identities such as \(2+3=5\) and
\(2\cdot 2=4\).
\end{summarybox}
